\begin{center}\begin{tikzpicture}[font=\small,>=stealth,line join=round]
\draw[gray!40,rounded corners=3pt] (0,0) rectangle (14.2,8.3);\node at (7.1,7.700000000000001) {只存在右半边：表达式不能替代定义域检查};
\begin{scope}[shift={(4.1,3.183333333333333)},x=1.0cm,y=1.5333333333333332cm]\draw[->,gray] (-3,0)--(3,0) node[right]{$x$};\draw[->,gray] (0,-1)--(0,2) node[above]{$y$};\draw[blue,thick,solid,domain=0:2.6,samples=180] plot(\x,{sqrt(\x)});\fill[orange!85!black] (1,1) circle(2.2pt);\end{scope}\node[anchor=west,align=left,text width=5.0cm] at (8.4,4.0) {$y=\sqrt{x}$\\[10pt]定义域 $[0,+\infty)$\\[10pt]$1$ 能取，$-1$ 不能取\\[10pt]既不是奇函数，也不是偶函数};
\node[text width=12.6cm,align=center] at (7.1,.65) {负半轴没有对应的函数值；左右配对是使用奇偶性规则的第一步。};
\end{tikzpicture}\end{center}
